\documentclass[a4paper,11pt]{article}

\usepackage[utf8]{inputenc}
\usepackage[T1]{fontenc}
\usepackage[ngerman]{babel}
\usepackage{amsmath,amssymb}
\usepackage{geometry}
\geometry{margin=2.0cm}
\usepackage{enumitem}
\usepackage{fancyhdr}
\usepackage{xcolor}
\usepackage{tikz}
\usepackage{titlesec}

%================= Dark-Mode Farben (Catppuccin Mocha) =================
\definecolor{base}{HTML}{1E1E2E}
\definecolor{fg}{HTML}{CDD6F4}
\definecolor{accent}{HTML}{89B4FA}
\definecolor{accent2}{HTML}{F5C2E7}
\definecolor{gridcolor}{HTML}{45475A}

\pagecolor{base}

\newcommand{\transp}{^{\mathsf{T}}}
\newcommand{\diag}{\operatorname{diag}}
\newcommand{\tridiag}{\operatorname{tridiag}}

\newcommand{\schwierig}[1]{\hfill\textnormal{\color{accent2}\itshape (#1)}}

% Karierte Antwortfläche -- #1 = Höhe (z. B. 8cm).
\newcommand{\karo}[1]{\par\vspace{0.3cm}\noindent
  \begin{tikzpicture}
    \draw[step=5mm, gridcolor, very thin] (0,0) grid (\linewidth,#1);
  \end{tikzpicture}\par\vspace{0.3cm}}

\titleformat{\section}{\color{accent}\normalfont\Large\bfseries}{}{0pt}{}

\pagestyle{fancy}
\fancyhf{}
\lhead{\color{fg}\small Finite Differenzen}
\rhead{\color{fg}\small FLA \& Analytische Geometrie}
\cfoot{\color{fg}\thepage}
\renewcommand{\headrule}{{\color{gridcolor}\hrule height 0.4pt}}

\setlength{\parindent}{0pt}

\begin{document}
\color{fg}

\thispagestyle{fancy}

\begin{center}
  {\color{accent}\LARGE\bfseries Übungsaufgaben: Finite Differenzen}\\[0.4em]
  {\large Fortgeschrittene Lineare Algebra und Analytische Geometrie}\\[0.3em]
  {\color{gridcolor}\rule{\linewidth}{0.4pt}}
\end{center}

\vspace{0.4em}

Sechs Übungsaufgaben \emph{nur} zum Thema Finite Differenzen -- genau im Stil der Klausur
(Randwertproblem $-u'' = f$ mit $u(0)=u(1)=0$, zentraler Differenzenquotient, tridiagonales
Gleichungssystem $Au=b$ bzw.\ $Tu = h^2 f$). Die Aufgaben werden von Aufgabe~1 nach Aufgabe~6
schwerer. Lösungen in \texttt{results/aufgaben/loesungen\_04.pdf}.

\medskip
\textbf{Nützliche Formeln:}
\[
\text{2.\ Ableitung: } u''(x_i) \approx \frac{u_{i-1} - 2u_i + u_{i+1}}{h^2},
\qquad
-u'' = f \ \Rightarrow\ -u_{i-1} + 2u_i - u_{i+1} = h^2 f(x_i).
\]
$T = \tridiag(2; -1)$ ist symmetrisch, positiv definit, tridiagonal und dünn besetzt.

% ================================================================
\clearpage
\section*{Aufgabe 1 -- Differenzenquotienten \schwierig{leicht}}

Sei $f(x) = x^3$. Approximieren Sie an der Stelle $x_0 = 1$ mit Schrittweite $h = 0{,}1$:
\begin{enumerate}[label=\alph*)]
  \item die erste Ableitung $f'(1)$ mit dem \textbf{Vorwärts-}, \textbf{Rückwärts-} und
        \textbf{zentralen} Differenzenquotienten;
  \item die zweite Ableitung $f''(1)$ mit dem zentralen Differenzenquotienten 2.\ Ordnung
        $\dfrac{f(x_0-h) - 2f(x_0) + f(x_0+h)}{h^2}$;
  \item Vergleichen Sie die Näherungen mit den exakten Werten $f'(1)$ und $f''(1)$. Welcher
        Quotient ist am genauesten?
\end{enumerate}

\karo{9cm}

% ================================================================
\clearpage
\section*{Aufgabe 2 -- Gleichungssystem aufstellen \schwierig{mittel}}

Gegeben sei das Randwertproblem
\[
-u''(x) = 8, \qquad u(0) = u(1) = 0,
\]
mit den inneren Gitterpunkten $x_1 = 0{,}25,\ x_2 = 0{,}5,\ x_3 = 0{,}75$.
\begin{enumerate}[label=\alph*)]
  \item Bestimmen Sie die Schrittweite $h$ und leiten Sie mit dem zentralen Differenzenquotienten
        das lineare Gleichungssystem $A u = b$ her.
  \item Nennen Sie vier strukturelle Eigenschaften der Matrix $A$.
\end{enumerate}

\karo{10cm}

% ================================================================
\clearpage
\section*{Aufgabe 3 -- Gleichungssystem lösen \schwierig{mittel}}

Lösen Sie das in Aufgabe 2 hergeleitete System
\[
\begin{pmatrix} 32 & -16 & 0 \\ -16 & 32 & -16 \\ 0 & -16 & 32 \end{pmatrix}
\begin{pmatrix} u_1 \\ u_2 \\ u_3 \end{pmatrix}
= \begin{pmatrix} 8 \\ 8 \\ 8 \end{pmatrix}.
\]
\begin{enumerate}[label=\alph*)]
  \item Lösen Sie es (Tipp: aus der Symmetrie folgt $u_1 = u_3$).
  \item Die exakte Lösung des Randwertproblems ist $u(x) = 4x(1-x)$. Vergleichen Sie Ihre
        Näherungswerte mit $u(0{,}25), u(0{,}5), u(0{,}75)$.
\end{enumerate}

\karo{10cm}

% ================================================================
\clearpage
\section*{Aufgabe 4 -- Feineres Gitter \schwierig{mittel--schwer}}

Gegeben sei
\[
-u''(x) = 2, \qquad u(0) = u(1) = 0,
\]
diesmal mit \textbf{vier} inneren Gitterpunkten $x_1 = 0{,}2,\ x_2 = 0{,}4,\ x_3 = 0{,}6,\ x_4 = 0{,}8$.
\begin{enumerate}[label=\alph*)]
  \item Geben Sie $h$ an und stellen Sie das $4\times 4$-System $T u = h^2 f$ mit $T = \tridiag(2;-1)$ auf.
  \item Lösen Sie das System (Tipp: $u_1 = u_4$ und $u_2 = u_3$).
\end{enumerate}

\karo{11cm}

% ================================================================
\clearpage
\section*{Aufgabe 5 -- Nichtkonstante rechte Seite \schwierig{schwer}}

Gegeben sei
\[
-u''(x) = x, \qquad u(0) = u(1) = 0,
\]
mit den inneren Gitterpunkten $x_1 = 0{,}25,\ x_2 = 0{,}5,\ x_3 = 0{,}75$.
\begin{enumerate}[label=\alph*)]
  \item Stellen Sie das System $T u = h^2 f$ auf. \textbf{Achtung:} die rechte Seite ist hier
        \emph{nicht} konstant -- $f(x_i) = x_i$ ändert sich pro Gitterpunkt.
  \item Lösen Sie das System.
\end{enumerate}

\karo{11cm}

% ================================================================
\clearpage
\section*{Aufgabe 6 -- Verständnisfragen \schwierig{Theorie}}

Beantworten Sie kurz und mit Begründung:
\begin{enumerate}[label=\alph*)]
  \item Warum koppelt jede Zeile des Systems nur $u_{i-1}, u_i, u_{i+1}$ -- und warum ist $A$
        deshalb \textbf{tridiagonal}?
  \item $A$ ist symmetrisch positiv definit. Welche Konsequenzen hat das für die Lösbarkeit und
        für die Wahl des Lösungsverfahrens?
  \item Wie sieht die rechte Seite $b$ aus, wenn $f$ \emph{nicht} konstant ist?
  \item Der zentrale Differenzenquotient 2.\ Ordnung hat Fehlerordnung $\mathcal{O}(h^2)$. Was
        passiert mit dem Diskretisierungsfehler ungefähr, wenn man $h$ halbiert?
  \item Warum tauchen die Randwerte $u(0)=u(1)=0$ in $b$ \emph{nicht} explizit auf?
\end{enumerate}

\karo{10cm}

\end{document}
